\section{Derivation of the fine-structure constant}
\label{sec:alpha}

In this development, the fine-structure constant receives an arithmetic
determination combining the three regional coordinates and the integer
dodecaphasic register. The operation combines their blocks with a fixed orientation and
transports the quotients of positional normalization. The resulting
parameter expresses the compatibility of that oriented difference with
the extension of the cylinders. Its value and the incidence data
used in its construction will remain related throughout the
exceptional development.

The internal order matters. The coordinates \(\pi\),
\(e\) and \(\varphi\), already constructed in the preceding section, are prior coordinates
of the same arithmetic state with trajectory memory, and may legitimately act as
inputs to a subsequent reader. They are not conventional constants
injected to select alpha. The correlated generation of the four coordinates shares an origin and
coinductive compatibility, without requiring all four
components to be evaluated in a single step.

The two routes express different mathematical functions of the same parameter.
The first performs an integer closure by Euclidean division in base
\(1000\). The second studies the analytic equilibrium of vacancies through
the resolvent \(E_{21/22}\) and its graded extensions. One route
determines blocks and quotients; the other determines a root through an
equation and a domain of uniqueness. Their identification is then formulated
in the common cylinders at every depth. The exposition distinguishes
these three results: positional construction, analytic determination
and compatibility of the two.

\subsection{First route: positional determination by integer closure}
\subsubsection{Domain and orientation of the integer closure}
\label{subsec:alpha-dominio}

Let \(P_j,E_j,\Phi_j\in\{0,\ldots,999\}\) be the canonical blocks of
the three regional coordinates, and extend the dodecaphasic register by
\[
 K_j=K_{1+((j-1)\bmod12)}.
\]
The channel orientation is \((1,1,-1)\). The contribution of the dodecaphasic register is
subtracted according to the dodecaphasic closure already typed. For \(B=1000\), the
local operation distinguishes the signed balance
\[
 Z_j:=P_j+E_j-\Phi_j-K_j
\]
from the sum incorporating the incoming quotient. Normalization is
\begin{equation}
 s_j=Z_j+c_{j+1},\qquad
 A_j=s_j-B\left\lfloor\frac{s_j}{B}\right\rfloor,
 \qquad c_j=\left\lfloor\frac{s_j}{B}\right\rfloor.
 \label{eq:alpha-recurrencia}
\end{equation}
Thus, \(A_j\in\{0,\ldots,B-1\}\), even when \(s_j\) is
negative. The carry entering from the right belongs to the domain of
the operation; it cannot be omitted because the table is printed from left
to right.

\begin{definition}[Finite-depth closure]
\label{def:alpha-finita}
For a depth \(N\) and an integer boundary \(t\), the map
\(\mathcal C_N\) receives
\((P_{1:N},E_{1:N},\Phi_{1:N},K_{1:N},c_{N+1}=t)\) and applies
\eqref{eq:alpha-recurrencia} in the order \(N,N-1,\ldots,1\). Its output
is the pair of words
\[
 (A_{1:N},c_{1:N+1}).
\]
\end{definition}

Euclidean division leaves no residual choice of sign: the remainder is
fixed in \(\{0,\ldots,999\}\) and the quotient is determined. For
example, \(-431=569-1000\) and \(-1200=800-2\cdot1000\). The
negative carries of the closure are not anomalies; they are the integer
information required to maintain the oriented equality.

\subsubsection{Dodecaphasic window and carry conditions}
\label{subsec:alpha-ventana}

The first twelve triads obtained are
\begin{align*}
 P={}&(141,592,653,589,793,238,462,643,383,279,502,884),\\
 E={}&(718,281,828,459,045,235,360,287,471,352,662,497),\\
 \Phi={}&(618,033,988,749,894,848,204,586,834,365,638,117).
\end{align*}
In the closure channel, the five regions determining the coordinate
\(P\) retain their distinct antecedents. Sharing this reading does not
identify them as regions or as histories.

With boundary \(c_{13}=0\), the recurrence produces the following table.
It shows each input, the incoming carry, the total before normalization,
the resulting block and the outgoing carry. The table is a reproducible integer
calculation, not a list of blocks selected for their resemblance to
a target.

\begin{table}[htbp]
\centering
\small
\setlength{\tabcolsep}{3.5pt}
\begin{tabular}{r|rrrr|rr|rr}
\hline
\(j\)&\(P_j\)&\(E_j\)&\(\Phi_j\)&\(K_j\)&\(c_{j+1}\)&\(s_j\)&\(A_j\)&\(c_j\)\\
\hline
1&141&718&618&234&0&7&007&0\\
2&592&281&033&543&0&297&297&0\\
3&653&828&988&140&\(-1\)&352&352&0\\
4&589&459&749&729&\(-1\)&\(-431\)&569&\(-1\)\\
5&793&045&894&659&\(-2\)&\(-717\)&283&\(-1\)\\
6&238&235&848&824&\(-1\)&\(-1200\)&800&\(-2\)\\
7&462&360&204&621&0&\(-3\)&997&\(-1\)\\
8&643&287&586&058&\(-1\)&285&285&0\\
9&383&471&834&914&\(-1\)&\(-895\)&105&\(-1\)\\
10&279&352&365&794&0&\(-528\)&472&\(-1\)\\
11&502&662&638&146&0&380&380&0\\
12&884&497&117&601&0&663&663&0\\
\hline
\end{tabular}
\caption{Initial closure with a zero right boundary. The dependence propagates from the last row to the first.}
\label{tab:alpha-ventana}
\end{table}

The carry word is
\[
 (c_1,\ldots,c_{13})=(0,0,0,-1,-1,-2,-1,0,-1,-1,0,0,0),
\]
and the reading of the window is the rational number
\begin{equation}
 \alpha_{12}^{(0)}
 =\frac{7297352569283800997285105472380663}{10^{36}}
 =0.007297352569283800997285105472380663.
 \label{eq:alpha-ventana-cero}
\end{equation}
The superscript recalls the chosen boundary. This window is exact in its
finite domain; it is not simply identified with the truncation of the
infinite section. The difference between the two readings becomes visible when
the carry from the remote tail is transported.

\subsubsection{Telescoping and reversibility in a boundary fiber}
\label{subsec:alpha-telescopia}

The local form of the equality is
\begin{equation}
 P_j+E_j-\Phi_j-K_j-A_j=Bc_j-c_{j+1}.
 \label{eq:alpha-identidad-local}
\end{equation}
The carry term links adjacent positions. Removing it
coordinate by coordinate yields a false equality. Only an
evaluation preserving its endpoints allows it to telescope.

\begin{proposition}[Integer identity at arbitrary depth]
\label{prop:alpha-telescopia}
For \(\iota_N(z)=\sum_{j=1}^{N}z_jB^{N-j}\), every execution of
\(\mathcal C_N\) satisfies
\begin{equation}
 \iota_N(P)+\iota_N(E)-\iota_N(\Phi)-\iota_N(K)-\iota_N(A)
 =B^Nc_1-c_{N+1}.
 \label{eq:alpha-telescopia}
\end{equation}
\end{proposition}

\begin{proof}
Multiplying \eqref{eq:alpha-identidad-local} by \(B^{N-j}\) and summing
produces on the right-hand side
\[
 \sum_{j=1}^{N}(Bc_j-c_{j+1})B^{N-j}=B^Nc_1-c_{N+1}.
\]
Each interior carry appears twice with opposite signs; only
the two endpoints survive.
\end{proof}

For Table~\ref{tab:alpha-ventana}, both endpoints are zero. If
\(p_{12},e_{12},f_{12}\) are the three truncated fractional sums,
\[
 \alpha_{12}^{(0)}=p_{12}+e_{12}-f_{12}-\kappa_{36}.
\]
The dodecaphasic register constant appearing here is \(\kappa_{36}\), not
\(\kappa_{\mathrm{per}}\). The domain of this identity contains
truncations and a finite boundary; its form does not authorize removing both
conditions when passing to the complete section.

The local inversion is also exact:
\[
 K_j=P_j+E_j-\Phi_j+c_{j+1}-A_j-Bc_j.
\]
Given \(P,E,\Phi\), the blocks \(A\) and the complete carries,
it recovers the dodecaphasic register. Concatenation also recovers
\[
 \iota_N(K)=\iota_N(P)+\iota_N(E)-\iota_N(\Phi)-\iota_N(A)
             -B^Nc_1+c_{N+1}.
\]
This reversibility is a property of the fiber with fixed data and
boundaries. It must not be turned into a new causal recipe that starts from
alpha and then manufactures the register. The direction of construction remains
signed register, dodecaphasic register, closure and output.

% BEGIN R27_MG05_CONTROLES
% MG05b. After inversion in the boundary fibre and before
% subsec:alpha-completa; requires tab:alpha-ventana and eq:alpha-telescopia.
\subsubsection{Controls of orientation, boundary, and register position}
\label{mg05:controles}

Integer closure distinguishes three operational dependencies:
the direction of carry propagation, the boundary condition, and the
position of the dodecaphasic register. Their comparison uses the same
twelve-block window as Table~\ref{tab:alpha-ventana}. The regional
words \(P,E,\Phi\) and the already generated register \(K\) are
retained except for the modification identified in each case. The
common reference is \(\mathcal C_{12}\) with \(c_{13}=0\), whose
output is \eqref{eq:alpha-ventana-cero}.

Prospective production of the signed register and its inversion
\(K^{(r)}=H_4U^{(r)}/4\) precede these controls. The normalization
of Lemma~\ref{lem:k-hadamard} is retained, in particular
\(H_4^2=4I_4\) and \(\det H_4=-16\). Inversion determines the
register once \(U\) has been produced; its operational selection
remains in the construction preceding closure. The following
variations of \(K\) examine the response of an already defined
reader.

\begin{proposition}[Separation of the dependencies of closure]
\label{mg05:dependencias}
For the data in Table~\ref{tab:alpha-ventana}, the following four
modifications have distinct effects:
\begin{enumerate}
\item Left-to-right quotient propagation, initialized with zero
quotient, changes the third block from \(352\) to \(353\).
\item The condition \(c_{13}=1\), with the original orientation,
changes the final block from \(663\) to \(664\) and preserves the
preceding eleven blocks.
\item The rotation \(K'_j=K_{j+1}\), with \(K_{13}=K_1\), changes
the output word while retaining \(c_{13}=0\).
\item The perturbation \(K'_1=K_1+1=235\), with the other eleven
coordinates fixed, changes the first block from \(007\) to \(006\)
and preserves the remaining blocks.
\end{enumerate}
\end{proposition}
\begin{proof}
For the first case, define the initial quotient \(b_0=0\) and
apply, in the order \(j=1,\ldots,12\),
\[
 \widetilde s_j=P_j+E_j-\Phi_j-K_j+b_{j-1},\qquad
 \widetilde A_j=[\widetilde s_j]_{1000},\qquad
 b_j=\frac{\widetilde s_j-\widetilde A_j}{1000}.
\]
This orientation produces, in two consecutive halves,
\begin{align*}
 \widetilde A_{1:6}&=(007,297,353,570,284,800),\\
 \widetilde A_{7:12}&=(995,285,106,471,379,663).
\end{align*}
The difference already appears at \(j=3\), where the quotient
required by the original closure comes from the fourth position.

For the second case, the final sum changes from \(663\) to \(664\).
Both have zero Euclidean quotient; hence \(c_{12}=0\) in both
cases and the recurrence agrees from \(j=11\) to \(j=1\).
The resulting word is
\begin{align*}
 A^{(1)}_{1:6}&=(007,297,352,569,283,800),\\
 A^{(1)}_{7:12}&=(997,285,105,472,380,664).
\end{align*}
Identity \eqref{eq:alpha-telescopia} records exactly the unit
variation of the right endpoint.

The rotation in the third case, evaluated with the same zero boundary
and the original recurrence, gives
\begin{align*}
 A^{\mathrm{rot}}_{1:6}&=(698,699,763,639,119,004),\\
 A^{\mathrm{rot}}_{7:12}&=(559,429,226,119,926,030).
\end{align*}
In this case the left carry is \(c_1=-1\); the right endpoint
remains zero. Retaining both endpoints in
\eqref{eq:alpha-telescopia} preserves the integer equality as the
twelve coefficients change position.

Finally, perturbing \(K_1\) leaves the entire recurrence from
\(j=12\) to \(j=2\) unchanged. At \(j=1\), the sum decreases
from seven to six, with zero quotient in both cases. This gives
\begin{align*}
 A^{\mathrm{pert}}_{1:6}&=(006,297,352,569,283,800),\\
 A^{\mathrm{pert}}_{7:12}&=(997,285,105,472,380,663).
\end{align*}
\end{proof}

The four controls separate transport direction, boundary fibre,
coefficient position, and local response. Their interpretation
retains the finite domain \(N=12\). Tail transport and normalization
of the complete section are developed next, with the compatibilities
required by passage to arbitrary depth.

% END R27_MG05_CONTROLES

\subsubsection{Extension and canonical normalization of the complete section}
\label{subsec:alpha-completa}

The extension preserves the periodic dodecaphasic register and receives the complete regional
sequences. Let
\[
 p=\sum_{j\geq1}P_jB^{-j},\qquad
 e_0=\sum_{j\geq1}E_jB^{-j},\qquad
 f=\sum_{j\geq1}\Phi_jB^{-j}.
\]
Each series is a realization of the corresponding internal character.
The integer parts already obtained are 3, 2 and 1, so that
\(p=\pi-3\), \(e_0=e-2\) and
\(f=\varphi-1\).

\begin{proposition}[Existence and uniqueness of the complete normalization]
\label{prop:alpha-normalizacion}
The series
\begin{equation}
 y=p+e_0-f-\kappa_{\mathrm{per}}
   =\sum_{j\geq1}(P_j+E_j-\Phi_j-K_j)B^{-j}
 \label{eq:alpha-serie}
\end{equation}
converges absolutely. In the domain of the blocks obtained,
\(0<y<10^{-2}\). There exists a unique canonical expansion \((A_j)\) of
\(y\), with no tail eventually constant at 999, and a unique bounded integer
sequence \((c_j)\) satisfying \eqref{eq:alpha-identidad-local} with
\(c_1=0\). These carries belong to the invariant set
\(\mathcal C=\{-2,-1,0,1,2\}\).
\end{proposition}

\begin{proof}
The coefficients of the series are bounded in absolute value by \(2(B-1)\),
so a geometric series dominates it. If
\(Z_j=P_j+E_j-\Phi_j-K_j\) and
\(R_j(Z)=\sum_{r\geq0}Z_{j+r}B^{-r-1}\), each tail is the difference
of two sums of two canonical tails and belongs to \((-2,2)\).
The first coefficient is
\(Z_1=141+718-618-234=7\); therefore
\(y=(7+R_2(Z))/1000\in(0.005,0.009)\).

Let \((A_j)\) be the canonical expansion of \(y\). Define
\[
 c_j=R_j(Z)-R_j(A).
\]
Multiplying the equality of the two series by \(B^{j-1}\) and separating
their first \(j-1\) terms gives
\[
 c_j=\sum_{r=1}^{j-1}(A_r-Z_r)B^{j-1-r}\in\mathbb Z,
\]
and \(c_1=0\). Moreover, \(R_j(A)\in[0,1)\), hence
\(-3<c_j<2\), which places all carries in \(\mathcal C\).
The shift identities for the tails give
\(Bc_j=Z_j-A_j+c_{j+1}\), which is the required recurrence.

Conversely, any bounded solution with canonical output and
\(c_1=0\) telescopes in the limit to the same series \(y\). The canonical
expansion is unique and, once fixed, the tail formula determines
all carries. Finally, if a finite execution is applied with
\(c_{j+1}\in\mathcal C\), then \(-2000\leq s_j\leq2000\),
so its Euclidean quotient again belongs to \(\mathcal C\).
\end{proof}

The proposition makes the normalization of the subsequent reader explicit: it does not
replace the generation of the three sequences with a conventional series
introduced as a datum. Once those sequences have been produced by the coinductive
state, their combination and the normalization boundary are
determined without consulting alpha.

The finite procedure preserved in the corpus propagates the five
boundaries of \(\mathcal C\) and retains only the prefix common to
all of them. If two depths have coalesced before a boundary,
the deterministic right-to-left recurrence gives the same blocks
to its left. This is compatibility with truncation. A
boundary is not chosen because it yields favorable digits, nor is
compatibility at every depth deduced from observing a long run.

The initial comparison shows why this precaution matters. The
stabilized extension gives \(c_{13}^{\infty}=-1\), whereas
the isolated window imposed \(c_{13}=0\). Thus,
\[
 884+497-117-601-1=662
\]
in the last block of the first turn. The first eleven triads do not
change, but the extension begins
\[
 (007,297,352,569,283,800,997,285,105,472,380,662,999,\ldots).
\]
An occurrence of the block 999 does not violate the canonical convention: what
is excluded is a tail eventually constant at 999. The window with a
zero boundary ends in 663; the truncation of the complete section
ends in 662. Both statements describe different and
exact operations. They must not be merged into a single projection notation.

\subsubsection{Operational definition and complete additive equation}
\label{subsec:alpha-definicion}

\begin{definition}[Alpha coordinate of the closure]
\label{def:alpha-operativa}
The integer route determines the coordinate
\begin{equation}
 \alpha_A=\sum_{j\geq1}A_jB^{-j}
 =p+e_0-f-\kappa_{\mathrm{per}},
 \label{eq:alpha-via-a}
\end{equation}
where \((A_j,c_j)\) is the compatible canonical normalization of
\eqref{eq:alpha-recurrencia}. In the scalar realization of the
HMT coordinates,
\begin{equation}
 \boxed{\alpha_A=
 \pi+e-\varphi
 -4-\kappa_{\mathrm{per}}.}
 \label{eq:alpha-identidad-completa}
\end{equation}
\end{definition}

The integer four is not calibrated: it is
\(3+2-1\), the combination of the integer parts of the three coordinates.
The base one thousand comes from the cubic chart, and the period twelve from the dodecaphasic register. The
signs are those of the oriented channel and the closure, and the remote boundary is
the one imposed by the compatible section. These are the effective origins of
the coefficients of the equation.

At this level, alpha may be described as the scalar coordinate of the
closure defect between the three regional realizations and the rational
dodecaphasic return, calculated with compatible carry. The term
\emph{defect} denotes a determined algebraic difference, not an
arbitrary loss or an imperfection of the system. This description does not
suppress the genealogy: without the regions, the signed register, the
orientation and the carry law, the scalar equality would not be a
sufficient exposition of the construction.

Three objects that alpha notation may accompany in subsequent readers
must also remain separate: \(A\) is a word of
blocks; the negative pre-carry support is obtained from the sign of the
totals before normalization; and an exceptional radial vector belongs to
an incidence codomain. None of them is interchangeable with the
real value \(\alpha_A\).

\paragraph{Signed balance and sum with an incoming quotient.}
\label{ampl-alfa-z-s-aclaracion}
Positional division consists of two consecutive stages. With
\(K_j^{\mathrm{per}}=K_{1+((j-1)\bmod12)}\), which denotes the
periodic extension also denoted by \(K_j\) in
\eqref{eq:alpha-recurrencia}, the balance before quotient
transport is
\[
 Z_j:=P_j+E_j-\Phi_j-K_j^{\mathrm{per}}.
\]
The variable $s_j$ in \eqref{eq:alpha-recurrencia} denotes the sum that already
includes the incoming contribution, so that
\[
 s_j=Z_j+c_{j+1},\qquad
 A_j=Z_j+c_{j+1}-1000c_j.
\]
The incidence configuration
$H^-_\alpha$ uses the signs of the balance $Z_j$ in the initial window;
determining the digits also uses the linking quotients
$c_{j+1}$ and $c_j$. This distinction jointly preserves the signed
support and the normalization of the expansion.

% Focal draft addition, NOT a proof of equality of the two routes.
% Its positive proposition may be incorporated without altering the current formulas.
% The note PROCEDENCIA_Y_ALCANCE.md delimits the compatibility not reconstructed.
\paragraph{Exact dependence on the truncation boundary.}
\label{ampl-alfa-frontera-exacta}
The regional coordinates and the dodecaphasic register arise from the
common APP--TRIT--TPK generation. Once produced, their positional
composition has an explicit dependence on the reading
depth. Let $B=1000$,
\[
 Z_j=P_j+E_j-\Phi_j-K_j^{\mathrm{per}},\qquad
 y_N=\sum_{j=1}^{N}Z_jB^{-j},\qquad
 R_{N+1}(Z)=\sum_{r\geq0}Z_{N+1+r}B^{-r-1}.
\]
The complete series in \eqref{eq:alpha-serie} satisfies exactly
\[
 \alpha_A=y_N+B^{-N}R_{N+1}(Z),\qquad
 -2<R_{N+1}(Z)<2.
\]
This identity combines the contribution of the window with that of its
extension, both expressed in the same coordinates.

\begin{proposition}[Boundary transport and canonical cylinder]
\label{ampl-alfa-prop-frontera}
Let $A_j^{(N,t)}$ be the output of finite positional division with
$c_{N+1}=t\in\mathcal C=\{-2,-1,0,1,2\}$, and let
$a_N^{(t)}=\sum_{j=1}^{N}A_j^{(N,t)}B^{-j}$. For the registers of the
domain under consideration, we have
\begin{equation}
 \alpha_A-a_N^{(t)}
 =B^{-N}\bigl(R_{N+1}(Z)-t\bigr),
 \qquad
 |\alpha_A-a_N^{(t)}|<4B^{-N}.
 \label{ampl-alfa-eq-frontera}
\end{equation}
The quotient of the complete extension is
$c_{N+1}^{\infty}=\lfloor R_{N+1}(Z)\rfloor$. With this boundary,
\[
 0\leq\alpha_A-a_N^{(c_{N+1}^{\infty})}<B^{-N},
\]
so the resulting window is the canonical prefix of order $N$.
\end{proposition}
\begin{proof}
Telescoping \eqref{eq:alpha-identidad-local} gives
$a_N^{(t)}=y_N+tB^{-N}-c_1$. Since $Z_1=7$, each incoming state
belongs to $\mathcal C$ and the first dividend lies between $5$ and $9$,
we have $c_1=0$. Substituting this equality into the decomposition of the
series proves \eqref{ampl-alfa-eq-frontera}. The bound follows from
$R_{N+1}(Z)\in(-2,2)$ and $t\in\mathcal C$.

For the complete canonical normalization,
$c_{N+1}^{\infty}=R_{N+1}(Z)-R_{N+1}(A)$ and
$R_{N+1}(A)\in[0,1)$. The integrality of the quotient implies
$c_{N+1}^{\infty}=\lfloor R_{N+1}(Z)\rfloor$. The error formula
then becomes
$\alpha_A-a_N^{(c_{N+1}^{\infty})}=B^{-N}R_{N+1}(A)$, proving
membership in the canonical cylinder.
\end{proof}

The proposition identifies the meaning of the right boundary: it is the
integer part of the signed tail expressed at the scale of the next
position. The calculation over the five boundaries preserves this information
as a finite collection of possibilities until the blocks to the
left stabilize. The quantitative contraction arises from
depth and tail transport; its object is the normalization of
the coordinates already generated.

\paragraph{Graded domain of the analytic resolvent.}
\label{ampl-alfa-dominio-resolvente}
The analytic composition displayed in
\S\ref{subsec:alpha-resolvente} acts on
$\operatorname{span}\{e_7,e_8,e_9\}$ and satisfies $Ne_9=0$. Its
evaluation produces exactly $T_{7:9}$; degree $9$ is the last
component of that domain. The identity $N^3=0$ explains why the
resolvent is expressed through three summands. To describe an
extension to higher degrees, one must also specify
the transport between successive graded domains and its action on the
initial component of each domain. Local nilpotence and
global compatibility thus play distinct mathematical roles.


\subsubsection{Constructive determination and role of the dodecaphasic register}
\label{alpha-ampl-dependencias}

The preceding definition clarifies what it means to determine alpha from
the structure. The domain of the operation contains four registers of
known provenance: the three regional sequences and the integer phase
register. Normalization also incorporates a boundary condition
compatible with extension. Once these data are fixed, Euclidean
division determines each remainder and its quotient. Proposition
\ref{prop:alpha-normalizacion} extends that determination to the complete
system. Independence from a target physical value is therefore
a property of the order of construction: the value compared at the
end does not participate in selecting any of the arguments of this
map.

This causal independence must be distinguished from algebraic independence
between numbers. The former is established by identifying the domain, the
operations and the provenance of the coefficients; the latter would be a
statement about polynomial relations. The result used here is the
former. The construction determines a coordinate through registers
produced by APP--TRIT--TPK and allows the relations it satisfies to be
checked subsequently. This order makes it possible to use a coordinate already
generated in a later operation without turning it into an external parameter.

The register $K$ has a more precise role than that of a numerical
correction added to a formula. Its twelve positions arise from the
reversible transformation of the signed register. The equation
$U=\mathcal H_{12}K$ preserves the possibility of returning to that register,
while the transverse differences and the total sum allow another
reconstruction. In the periodic evaluation, the order of the positions
determines the positional weight of each component. The same set of
components, arranged in a different order, generally produces a different rational
coordinate. Thus, the phase origin and orientation form part of the
definition of the number entering \eqref{eq:alpha-identidad-completa}.

The normalization of alpha composes this information with the regions. Its
local identity distinguishes two successive operations. First the oriented
balance of blocks is formed; then that balance is represented through
a canonical remainder and a quotient linking adjacent positions. The first
operation preserves the sign of the balance and allows a support to be
extracted subsequently. The second allows the prefixes of the real
coordinate to be formed. Both readings remain related because the state preserves the
data before and after normalization. The signed support and the
positional expansion are thus different realizations of an explicit
arithmetic composition.

The telescoping identity expresses the global content of this composition.
The interior quotients cancel in pairs when a prefix is evaluated,
while their endpoints remain in the equality. This cancellation is a
property of the weighted sum; it neither removes the quotients from the state nor
allows them to be suppressed before the sum is taken. The right boundary of a
window receives the contribution of the continuation. For this reason,
evaluating an isolated window and restricting an extended
section are distinct operations. The change of the final block from $663$ to
$662$, proved above, provides an explicit realization of this
difference without altering the blocks already stabilized to its left.

In this sense, bidirectionality has a concrete arithmetic
formulation. Generation provides compatible extensions;
normalization transmits the integer information of their boundary from
positions of lower weight towards those of higher weight. The equation preserves the
contribution of both directions through its indices and its boundary
conditions. Reversibility in the fixed fiber is checked by reconstructing
$K$ from the other registers and the complete quotients. This
inversion certifies the consistency of the original operation; the genealogical
direction remains the one that produces the register first and the alpha
coordinate afterwards.

The second determination has complementary content. In it, the
immediate object is a balance of vacancies, and the parameter appears as
its root in an isolating domain. The positional route provides a
canonical normalization of registers; the analytic route studies the vanishing
of an operator. Presenting both requires precisely this
difference of functions to be preserved. Their relations with the common state must be
expressed through the refinement maps and the comparison of
their complete realizations, as developed in the corresponding
subsection. Agreement of a finite approximation constitutes a
localized check; identification of the complete constructions
has the domain and scope of its own statement.

Finally, this organization explains the position of alpha in the article.
Its equation appears after regional generation and the construction
of $K$, because it uses both results. The electron section then uses
that coordinate as a factor in its composition. The exceptional
extension receives the registers and their supports, where incidence
information is preserved. These two continuations share the same
provenance and develop distinct mathematical functions: evaluation of
a central section and construction of a boundary configuration.
Developing each one allows the structure to be followed through the
equation, instead of restricting the comparison to its scalar value.


\subsection{Second route: analytic determination through vacancies}
\subsubsection{Transport invariants and graded resolvent}
\label{subsec:alpha-via-b}

The analytic route receives invariants of the oriented state, not the output
\(\alpha_A\). Its sixth-order truncation is
\begin{equation}
 P_6(x)=2\pi x-\frac74x^2
 +\frac{x^3}{2\pi}
 +\frac{x^4}{20}-\frac{2x^5}{21}-\frac{x^6}{46},
 \label{eq:alpha-p6}
\end{equation}
and the vacancy condition uses
\[
 \delta=\log_{10}\frac{10}{9}.
\]
The ratio \(10/9\) compares the decimal chart with the nine APP positions
of the nonadic construction. Its logarithm belongs to this subsequent
compactification reading; it is not a target value of alpha.

The genealogy of the coefficients must be preserved in full, but with its
exact logical strength. Recent developments bring together a marked signature
of the state. The synchronization matrix and the associated census are
\[
 D_{\mathrm{syn}}=\begin{pmatrix}85&112\\41&52\end{pmatrix},
 \qquad N_9=43.
\]
Its determinant yields
\[
 a=-\frac{\det D_{\mathrm{syn}}}{N_9}
   =\frac{172}{43}=4.
\]
The clock induced on the short gaps has secondary spectrum
\(\{6,7\}\); its upper endpoint gives \(b=7\). The tritic
cardinality \(m=3\) determines
\[
 t_*=mb=21,\qquad k_*=mb-1=20.
\]
The vacancy reading in the base-one-thousand chart gives
\[
 v_-=\lfloor1000\delta\rfloor=45,\qquad
 v_+=\lceil1000\delta\rceil=46.
\]
These integers are typed antecedents of the evaluation of the truncation. They are not
obtained from its root or from a metrological approximation to it.

The lower branch also coincides with the determinant of the APP block
\[
 B_c=\begin{pmatrix}7&2\\2&7\end{pmatrix},\qquad\det B_c=45.
\]
The transverse rank of the register in the preceding chapter gives
\(r_{\mathrm{dod}}=11\), and the phase half-corona provides
\(f_{\mathrm{ph}}=8\binom62=120\). In particular,
\(45r_{\mathrm{dod}}=495\).

With \(\omega=2\pi\), the signature
\[
 \mathcal I_9=(\omega,m,a,b,k_*,t_*,v_-,v_+,f_{\mathrm{ph}},r_{\mathrm{dod}})
\]
feeds the graded reader
\begin{align}
 \mathfrak E_9(\mathcal I_9)(x)
 ={}&\omega x-\frac ba x^2+\omega^{-1}x^3
 +\frac{x^4}{k_*}-\frac{2x^5}{t_*}-\frac{x^6}{v_+}\notag\\
 &-\frac{x^7}{f_{\mathrm{ph}}}+\frac{x^8}{v_-}
 +\frac{2x^9}{v_-r_{\mathrm{dod}}}.
 \label{eq:alpha-firma-coeficientes}
\end{align}
Substitution recovers the coefficients of \(P_6\) and its
ninth-order extension. The pair \(45/46\) crosses the division
between the two: the upper branch occupies degree six, and the lower branch reappears
in degrees eight and nine.

The factorization in \eqref{eq:alpha-firma-coeficientes} is formulated
for a state preserving origin, channel order and orientation.
Naturality with respect to morphisms preserving those data is not equivalent to
uniqueness of every reader after they have been forgotten. The domain of this
factorization is determined by the signature
\eqref{eq:alpha-firma-coeficientes}; the extension to all degrees
is constructed in Section~\ref{subsec:alpha-publicaciones-completas},
with its recurrence class and its dependence on the state's pre-coordinate
explicitly defined.

\subsubsection{Nilpotent resolvent in degrees 7, 8 and 9}
\label{subsec:alpha-resolvente}

The first extension admits a particularly explicit finite operator
construction. On the space with basis \(e_7,e_8,e_9\),
let
\[
 Ne_7=-\frac83e_8,\qquad Ne_8=\frac2{11}e_9,\qquad Ne_9=0,
 \qquad v_7=-\frac1{120}e_7.
\]
The symbol \(N\) here denotes a nilpotent operator, not a truncation
depth. Its domain and grading are fixed before the polynomial
is evaluated at a number.

\begin{lemma}[Exact nilpotent extension]
\label{lem:alpha-resolvente}
We have \(N^3=0\) and
\[
 (I-N)^{-1}v_7
 =-\frac{e_7}{120}+\frac{e_8}{45}+\frac{2e_9}{495}.
\]
The evaluation \(e_n\mapsto x^n\) produces
\begin{equation}
 T_{7:9}(x)=-\frac{x^7}{120}+\frac{x^8}{45}+\frac{2x^9}{495}.
 \label{eq:alpha-t79}
\end{equation}
\end{lemma}

\begin{proof}
The operator raises the degree and vanishes on the last basis vector,
so its cube is zero. The inverse is the finite resolvent
\(I+N+N^2\). Moreover,
\[
 Nv_7=\frac1{45}e_8,\qquad N^2v_7=\frac2{495}e_9.
\]
Summing the three terms and evaluating gives the formula.
\end{proof}

Here the denominators 120, 45 and 495 are not decimal adjustments: they arise
from the half-corona, the APP determinant and its product with the dodecaphasic
rank. The transfers \(-8/3\) and \(2/11\), together with the
sign and degree of the seed, belong to the typed reader. Preserving
only the symmetry of the cycle and replacing this operator with another symmetric
operator does not necessarily preserve the coefficients.

\subsubsection{Root of multiplicity one of the ninth-order truncation}
\label{subsec:alpha-jet}

Define
\[
 P_9=P_6+T_{7:9},\qquad E_{21/22}^{(9)}(x)=P_9(x)-\delta,
 \qquad I_\alpha=(0,10^{-2}).
\]
This finite operator has its own existence and uniqueness result.

\begin{proposition}[Root of multiplicity one of the truncation]
\label{prop:alpha-raiz-jet}
The operator \(E_{21/22}^{(9)}\) has a unique root
\(\xi_9\) in \(I_\alpha\), and that root is simple. The following rational isolation is obtained:
\begin{equation}
 \frac{729735256928380099}{10^{20}}
 <\xi_9<\frac{729735256928380100}{10^{20}}.
 \label{eq:alpha-intervalo-jet}
\end{equation}
\end{proposition}

\begin{proof}
The bounds on the internal coordinates
\(3.14<\pi<3.15\) and
\(0.045<\delta<0.046\) suffice for existence. At zero,
\(E_{21/22}^{(9)}(0)=-\delta<0\). At \(10^{-2}\), the linear
contribution minus the negative terms is greater than \(0.0626\), so the
operator's value is positive.

For \(0\leq x\leq10^{-2}\), discarding the positive terms of
the derivative gives
\[
 (E_{21/22}^{(9)})'(x)
 \geq2\pi-\frac72\cdot10^{-2}
 -\frac{10}{21}\cdot10^{-8}-\frac6{46}\cdot10^{-10}
 -\frac7{120}\cdot10^{-12}>6.24.
\]
Continuity gives existence; the derivative bound gives uniqueness and
simplicity. The narrower interval
\eqref{eq:alpha-intervalo-jet} results from a directed rational evaluation.
Let \(x_-\) and \(x_+\) be its endpoints. To compute
\(\delta=\log(10/9)/\log 10\), we use, for \(z>1\),
\[
 L_N(z)=2\sum_{j=0}^{N-1}\frac{y^{2j+1}}{2j+1},\qquad
 y=\frac{z-1}{z+1},\qquad
 0<\log z-L_N(z)<\frac{2y^{2N+1}}{(2N+1)(1-y^2)}.
\]
The remainder is bounded by the geometric series of positive terms.
This formula is applied to \(z=10/9\) and \(z=3\), with \(N=520\),
using \(\log 10=\log(10/9)+2\log3\). Evaluation of
\(\pi\) uses the generated thousand-digit regional cylinder, including
its width \(10^{-1000}\). Interval arithmetic in
\eqref{eq:alpha-firma-coeficientes} gives
\[
 -\frac{4559}{10^{23}}<E_{21/22}^{(9)}(x_-)< -\frac{4558}{10^{23}},
 \qquad
 \frac{1698}{10^{23}}<E_{21/22}^{(9)}(x_+)<\frac{1699}{10^{23}}.
\]
These signs and the proved monotonicity establish the isolation. The endpoints
are results of evaluating the operator; its coefficients arise from
the signature constructed beforehand.
\end{proof}

The result concerns \(\xi_9\). It does not authorize writing
\(\xi_9=\alpha_A\) as an exact equality of complete numbers.
The root of a truncation may share a prefix with the root of the
extended operator and differ at a later position. Simplicity
allows the motion of the root under perturbations to be controlled; it does not
eliminate those perturbations.

\subsection{Compatibility of the two determinations at every depth}
\label{subsec:alpha-limite-dos-vias}

\begin{remark}[Truncation, extension and compatibility]
The operations in degrees \(6\) to \(9\) determine a finite jet.
Positional normalization determines the complete coordinate through
its compatible sequence of carries. The next proposition expresses
the general criterion for identification through cylinders.
Section~\ref{subsec:alpha-publicaciones-completas} then constructs
a correlated analytic chart of all degrees, with a recurrence,
uniform convergence, uniqueness of the root and projective compatibility.
The construction receives the shared coordinate of the state and the declared
class of extensions; it thus preserves the distinction between an
analytic representation of that coordinate and an independent selection.
\end{remark}

The complete analytic route is organized into a collection of compatible graded
readers
\[
 E_{21/22}^{(6)}\longleftarrow E_{21/22}^{(9)}
 \longleftarrow E_{21/22}^{(12)}\longleftarrow\cdots.
\]
The symbol
\[
 E_{21/22}=\varprojlim_m E_{21/22}^{(6+3m)}
\]
denotes that complete reader, with its transports and truncations, not a
formal series freely chosen to make it pass through \(\alpha_A\).
The second route determines its unique positive root in \(I_\alpha\),
denoted by \(\alpha_B\), with the uniqueness and compatibility of the
complete reader proved in Section~\ref{subsec:alpha-publicaciones-completas}.

The proofs in Section~\ref{subsec:alpha-publicaciones-completas}
establish that, at each depth, both systems determine the same
surviving cylinder \(C_N^\alpha\), and that their diameters tend to
zero. It is useful to state the mathematical consequence of this
compatibility separately.

\begin{proposition}[Identification of the compatible limits]
\label{prop:alpha-dos-limites}
If the values of the integer and analytic routes belong to the
same cylinders \(C_N^\alpha\) for every \(N\), and
\(\operatorname{diam}C_N^\alpha\to0\), then
\[
 \alpha_A=\alpha_B=\alpha.
\]
\end{proposition}

\begin{proof}
For every \(N\), the distance between the two values is bounded
by \(\operatorname{diam}C_N^\alpha\). Letting \(N\) tend to infinity
gives zero distance. The existence of the compatible sections ensures
that this is not the empty intersection of a system of approximations.
\end{proof}

This proof identifies the consequence of the compatibility squares.
Their construction at arbitrary depth appears in
Section~\ref{subsec:alpha-publicaciones-completas}. Finite
evaluations check realizations of that construction; the recurrence
and its uniform bounds establish the passage to the limit.

The difference can be expressed algebraically. Let
\(F=\mathbb Q(\pi,\delta)\) and let
\(j_9:F[[x]]\to F[x]/(x^{10})\) be the truncation. Then
\[
 \ker j_9=x^{10}F[[x]].
\]
All the operators
\[
 E_h(x)=E_{21/22}^{(9)}(x)+x^{10}h(x)
\]
have the same ninth-order truncation. However, taking
\(h=0\) and \(h=1/729\),
\[
 E_{1/729}(\xi_9)=\frac{\xi_9^{10}}{729}>0,
\]
so they do not share that root. The observation does not allow a
tail foreign to HMT transport to be chosen; it demonstrates precisely why the generated
tail and its compatibility must be preserved. Nor does an arbitrary formal
series suffice to speak of analytic roots without its domain of
convergence or the realization of the cylinder system.
\subsection{Compatibility between the dodecaphasic and analytic representations of alpha}
\label{subsec:alpha-publicaciones-completas}

Compatibility is established at every depth. The reversible
dodecaphasic representation and the analytic representation of the alpha
coordinate
\(E_{21/22}\) are two internal representations with distinct codomains,
produced from the same enriched APP--TRIT--TPK state. The jet
\(E_{21/22}^{(9)}\) is an exact finite witness, but does not by itself determine
the tail. The completion constructed below is not deduced from the bare
jet: it receives the shared coordinate of the state and expresses its
analytic image through an explicit recurrence. It therefore does not constitute
an independent causal selection; it constitutes a typed analytic
representation of the same generated point.

This separation preserves the causal order
\[
 \mathrm{APP}\longrightarrow\mathrm{TRIT}\longrightarrow\mathrm{TPK}
 \longrightarrow\mathcal X_{\alpha}^{\mathrm{enr}}
 \longrightarrow (P,E,\Phi,U_{12}^{\mathrm{sgn}},K)
 \longrightarrow\alpha_{\mathrm{HMT}}.
\]
The coordinates \(P,E,\Phi\) and the register \(K\) are earlier outputs of
the same enriched state. Their use in the closure of alpha is therefore a
subsequent internal composition, not an injection of conventional
constants.

For \(r\in6\mathbb N\), \(r\ge36\), let
\(\widehat{\mathcal X}^{\mathrm{enr},\alpha}_r\) be the sector of truncated
states that extend the internal alpha seed fixed at \(R_{36}\). Write
\[
 x_r^\alpha=(\widetilde x_r,\chi_\alpha,
              \varepsilon_r^\alpha,N_r^\alpha),
\]
where \(\chi_\alpha\) is the internal character of the sector, not the expressed
scalar value. Each state preserves remainder, quotient, carry, sheet, orientation,
phase, route, cylinder, boundary, signature, register, memory, survival,
provenance and depth. If
\[
 d_r^\alpha=\operatorname{Dig}_{729}(\varepsilon_r^\alpha),\qquad
 \varepsilon_{r+6}^\alpha=729\varepsilon_r^\alpha-d_r^\alpha,
 \qquad N_{r+6}^\alpha=729N_r^\alpha+d_r^\alpha,
\]
the effective transition is
\[
 \mathcal U_r^\alpha=
 \operatorname{Upd}_r^\alpha\circ
 \operatorname{Tra}_r^\alpha\circ
 \operatorname{Sel}_r^\alpha,
 \qquad
 \operatorname{Ext}^{\alpha,\to}_{r+6,r}
 =\operatorname{graph}(\mathcal U_r^\alpha).
\]
The single-valuedness of these extensions, with the seed \(R_{36}\) fixed,
produces a unique compatible section: the phase returns while memory and
depth advance. The global domain of both determinations is
\begin{equation}
 \mathcal X_\alpha^{\mathrm{enr}}:=\widehat\Omega_\alpha
 :=\varprojlim_{r\ge36}
 \bigl(\widehat{\mathcal X}^{\mathrm{enr},\alpha}_r,
       \operatorname{Ext}^{\alpha,\to}_{r+6,r}\bigr).
 \label{eq:alpha-dominio-enriquecido-completo}
\end{equation}

In this section, \(\kappa_{\mathrm{per}}\) is the periodic reading \(\kappa_{\mathrm{ag}}\) of the register already constructed; both notations denote the same fraction in the fixed base-one-thousand chart.

\subsubsection{The dodecaphasic closure as a complete section}

For \(N\geq1\), set
\[
 \mathcal W_N:=\{0,\ldots,999\}^{N},
 \qquad
 \mathsf T_{\alpha,N}(\mathsf s):=(A_1,\ldots,A_N),
 \qquad \mathsf s\in\mathcal X_{\alpha}^{\mathrm{enr}},
\]
where the blocks \(A_j\) are obtained through
\eqref{eq:alpha-recurrencia} with the compatible boundary
\(c_{N+1}^{\infty}=\lfloor R_{N+1}(Z)\rfloor\). If \(N'\geq N\), let
\(\rho_{N',N}:\mathcal W_{N'}\to\mathcal W_N\) be the restriction to the
first \(N\) blocks.

\begin{theorem}[Complete dodecaphasic output and projective compatibility]
\label{thm:alpha-salida-dodecafasica-completa}
The collection \((\mathsf T_{\alpha,N})_{N\geq1}\) is compatible:
\begin{equation}
 \rho_{N',N}\circ\mathsf T_{\alpha,N'}
 =\mathsf T_{\alpha,N}
 \qquad(N'\geq N).
 \label{eq:alpha-via-a-proyectiva-exacta}
\end{equation}
Its canonical cylinders
\begin{equation}
 C_N^A(\mathsf s):=
 \left[
  \sum_{j=1}^{N}A_j1000^{-j},
  \sum_{j=1}^{N}A_j1000^{-j}+1000^{-N}
 \right)\cap I_\alpha
 \label{eq:alpha-cilindros-via-a-exactos}
\end{equation}
are nested, have diameter at most \(1000^{-N}\) and satisfy
\begin{equation}
 \bigcap_{N\geq1}C_N^A(\mathsf s)
 =\{\alpha_A(\mathsf s)\},
 \qquad
 \alpha_A
 =\pi_{\mathrm{HMT}}+e_{\mathrm{HMT}}-\varphi_{\mathrm{HMT}}
  -4-\kappa_{\mathrm{per}}.
 \label{eq:alpha-salida-a-completa}
\end{equation}
Consequently, the first route expresses a unique infinite section, which we
denote by \(\alpha_{\mathrm{HMT}}:=\alpha_A\).
\end{theorem}

\begin{proof}
Proposition~\ref{prop:alpha-normalizacion} provides the unique canonical
expansion and the unique bounded sequence of carries. The local identity
\eqref{eq:alpha-identidad-local} and its telescoping
\eqref{eq:alpha-telescopia} preserve the remote boundary; Proposition
\ref{ampl-alfa-prop-frontera} proves that restricting an extended execution
recovers exactly the preceding prefix. This gives
\eqref{eq:alpha-via-a-proyectiva-exacta}. Membership of the full sum
in each cylinder follows from the tail identity, and
\(\operatorname{diam}C_N^A\leq1000^{-N}\to0\); the intersection contains a
single point. Finally, \eqref{eq:alpha-identidad-completa} identifies that
point with the expression in \eqref{eq:alpha-salida-a-completa}. All its
coefficients and signs are fixed before the representation of alpha.
\end{proof}

The window \(\alpha_{12}^{(0)}\), closed with \(c_{13}=0\), remains an
exact finite calculation; it must not be confused with the restriction of the
complete section, whose compatible boundary satisfies \(c_{13}^{\infty}=-1\). In
particular,
\begin{equation}
 \rho_{36,12}(\alpha_{A,36})=\alpha_{A,12}
 \label{eq:alpha-rho-36-12-proyectiva}
\end{equation}
is an equality of compatible words, not an equality between words of
different lengths or an identification of the isolated window with the
limit.


% Correlated analytic representation of alpha: action, convergence and
% intervals.
\subsubsection{Analytic representation of the alpha coordinate}
\label{subsec:alpha-lector-completo-rev08}

The analytic extension is specified through all its coefficients.
Let \(\mathsf s\in\mathcal X_\alpha^{\mathrm{enr}}\). Before
decimal normalization, the same state already expresses the pre-coordinate
\begin{equation}
 a(\mathsf s):=p(\mathsf s)+e_0(\mathsf s)-f(\mathsf s)
                 -\kappa_{\mathrm{per}}(\mathsf s)\in I_\alpha,
 \label{eq:alpha-precoordenada-comun}
\end{equation}
with \(p=\pi_{\mathrm{HMT}}-3\),
\(e_0=e_{\mathrm{HMT}}-2\) and
\(f=\varphi_{\mathrm{HMT}}-1\). Equivalently,
\(a=\pi_{\mathrm{HMT}}+e_{\mathrm{HMT}}-
\varphi_{\mathrm{HMT}}-4-\kappa_{\mathrm{per}}\). No conventional
\(\alpha\) is introduced here: \(p,e_0,f\) are the three regional characters already
generated, and \(\kappa_{\mathrm{per}}\) is the rational reading of the terminal
register. The integer closure applies to
\eqref{eq:alpha-precoordenada-comun} the normalization in
\eqref{eq:alpha-recurrencia}; the chart below expresses a compatible
analytic representation on the same pre-coordinate. Its dependency
graph receives \((p,e_0,f,\kappa_{\mathrm{per}})\) directly: it does not
contain an edge \(\alpha_A\to\lambda\), although the equality proved below
allows both representation names to be identified subsequently.

Write \(q=1/729\), let \(F_{\mathsf s}\) be the ordered field generated by
the coordinates of \(\mathsf s\), and set
\(E_{9,\mathsf s}(X)=E_{21/22}^{(9)}(\mathsf s;X)\). We define the linking
coefficient
\begin{equation}
 \lambda(\mathsf s):=
 -\frac{E_{9,\mathsf s}(a(\mathsf s))
         \bigl(1-q\,a(\mathsf s)^3\bigr)}{a(\mathsf s)^{10}}.
 \label{eq:alpha-lambda-estado}
\end{equation}
All its arguments belong to the state before metrological
recognition. In particular, \(\lambda\) is not obtained by fitting a target
decimal string. Equation~\eqref{eq:alpha-lambda-estado} must be read with
its exact type: once the first representation of the state has produced
\(a(\mathsf s)\), it fixes the unique extension within the geometric class
defined below. It is not a second independent selection of
\(a(\mathsf s)\), but a correlated analytic representation of the same
enriched state.

Indeed, for \(\mu\in F_{\mathsf s}\) set
\[
 E_{\mathsf s,\mu}(X)
 :=E_{9,\mathsf s}(X)+\mu\frac{X^{10}}{1-qX^3}.
\]
Since \(0<a(\mathsf s)<10^{-2}\) and \(qa(\mathsf s)^3<1\), we have
\begin{equation}
 E_{\mathsf s,\mu}(a(\mathsf s))=0
 \quad\Longleftrightarrow\quad
 \mu=\lambda(\mathsf s).
 \label{eq:alpha-unicidad-lambda-clase-geometrica}
\end{equation}
Thus, the coefficient seed contains no freedom once the state
and the class of \(q\)-geometric extensions have been fixed. This uniqueness is relative to the
declared class; it does not assert that the ninth-order jet by itself selects a
tail among all possible analytic series.

The local action on blocks of three coefficients is
\begin{equation}
 \mathcal C_\alpha:F_{\mathsf s}^{3}\longrightarrow F_{\mathsf s}^{3},
 \qquad \mathcal C_\alpha(u,v,w)=q(u,v,w),
 \qquad c^{(0)}=(\lambda(\mathsf s),0,0).
 \label{eq:alpha-accion-coeficientes}
\end{equation}
Therefore, for every \(n\geq0\),
\begin{equation}
 b_{10+3n}=q^n\lambda(\mathsf s),\qquad
 b_{11+3n}=b_{12+3n}=0.
 \label{eq:alpha-recurrencia-coeficientes}
\end{equation}
This is the rule that computes each coefficient beyond degree nine. No
free parameter remains at each extension.

For \(M\geq0\), let
\begin{equation}
 E_{\mathsf s,M}(X):=E_{9,\mathsf s}(X)
 +\lambda(\mathsf s)\sum_{n=0}^{M}q^nX^{10+3n}.
 \label{eq:alpha-truncamientos-completos}
\end{equation}
The complete function is
\begin{equation}
 E_{\mathsf s,\infty}(X)
 :=E_{9,\mathsf s}(X)
 +\lambda(\mathsf s)\frac{X^{10}}{1-X^3/729}.
 \label{eq:alpha-funcion-completa-explicita}
\end{equation}

To formulate the projective system without implicit types, we fix
\(d_M=12+3M\) and define the jet fibration
\begin{equation}
 \mathfrak J_M:=
 \coprod_{\mathsf s\in\mathcal X_\alpha^{\mathrm{enr}}}
 \{\mathsf s\}\times F_{\mathsf s}[X]/(X^{d_M+1}).
 \label{eq:alpha-espacio-jets-rev08}
\end{equation}
If \(M'\ge M\), the restriction is
\begin{equation}
 \tau_{M',M}(\mathsf s,[P]_{d_{M'}})
 :=(\mathsf s,[P]_{d_M}).
 \label{eq:alpha-truncamiento-jets-rev08}
\end{equation}

\begin{proposition}[Naturality of the analytic recurrence]
\label{prop:alpha-naturalidad-recurrencia-explicita}
Let
\[
 \Gamma_M(\mathsf s)
 =\bigl(\mathsf s,[E_{\mathsf s,M}]_{d_M}\bigr).
\]
If \(M'\geq M\), the truncation \(\tau_{M',M}\) satisfies
\begin{equation}
 \tau_{M',M}\circ\Gamma_{M'}=\Gamma_M,
 \qquad
 \Gamma_{E21}(\mathsf s)=\varprojlim_M\Gamma_M(\mathsf s).
 \label{eq:alpha-naturalidad-explicita}
\end{equation}
The projection onto the three new degrees is precisely
\(\mathcal C_\alpha^{M+1}(c^{(0)})\).
\end{proposition}

\begin{proof}
Equation~\eqref{eq:alpha-recurrencia-coeficientes} fixes the block of
degrees \(10+3n,11+3n,12+3n\). Passing from \(M\) to \(M+1\) adds only
the next block; truncating it recovers the preceding polynomial literally.
Compatibility for \(M'\geq M\) is obtained by composition. Thus, the
inverse limit is not a section chosen from the bare jet, but the
unique orbit of \(c^{(0)}\) under \(\mathcal C_\alpha\).
\end{proof}

\subsubsection{Uniform convergence, simple root and rational intervals}
\label{subsec:alpha-cotas-completas-rev08}

The regional cylinders and the periodic register provide bounds through
rational interval arithmetic. For each coordinate
\[
 c\in\{\pi_{\mathrm{HMT}},e_{\mathrm{HMT}},\varphi_{\mathrm{HMT}}\},
\]
let \(P_c\) be its certified thousand-digit nonadic prefix. The datum used
is not the rational number \(P_c\) treated as though it were the complete value, but the
half-open cylinder and its outer closure
\[
 \mathcal C_c^{(1000)}=[P_c,P_c+\varepsilon),
 \qquad
 \overline{\mathcal C}_c^{(1000)}=[P_c,P_c+\varepsilon],
 \qquad \varepsilon=10^{-1000}.
\]
Directed arithmetic operates conservatively with these closures. If
\(P_\pi,P_e,P_\varphi\) denote the three prefixes, we obtain
\begin{equation}
\begin{aligned}
 a_-&:=P_\pi+P_e-(P_\varphi+\varepsilon)-4-
       \kappa_{\mathrm{per}},\\
 a_+&:=(P_\pi+\varepsilon)+(P_e+\varepsilon)-P_\varphi-4-
       \kappa_{\mathrm{per}},\\
 A_{1000}&:=[a_-,a_+],
 \qquad a(\mathsf s)\in[a_-,a_+]=A_{1000},
 \qquad \operatorname{diam}A_{1000}=3\varepsilon.
\end{aligned}
\label{eq:alpha-propagacion-colas-rev08}
\end{equation}
The natural interval extension of
\eqref{eq:alpha-lambda-estado} then transports \(A_{1000}\), the cylinder
of \(\pi_{\mathrm{HMT}}\) and the directed rational bound for \(\delta\) to a
closed interval \(\Lambda_{1000}\ni\lambda(\mathsf s)\); none of the three
tails is frozen. In particular,
\begin{equation}
\begin{aligned}
 \frac{7297352569283800997285105472380662}{10^{36}}
 &<a(\mathsf s)<
 \frac{7297352569283800997285105472380663}{10^{36}},\\
 -8.711\mathbin{\times}10^{-6}
 &<\lambda(\mathsf s)<-8.709\mathbin{\times}10^{-6}.
\end{aligned}
 \label{eq:alpha-cotas-precoordenada-lambda}
\end{equation}
These inequalities use certified prefixes with their rational tail
bound; they do not receive a table of alpha values. The package verifier reproduces
them from the three nonadic outputs and from the terminal chart.

Let
\[
 r:=\frac{10^{-6}}{729}<1.372\mathbin{\times}10^{-9}.
\]
For \(0\leq x\leq10^{-2}\), the tail beyond \(M\) satisfies
\begin{equation}
 \left|E_{\mathsf s,\infty}(x)-E_{\mathsf s,M}(x)\right|
 \leq
 8.711\mathbin{\times}10^{-26}\,\frac{r^{M+1}}{1-r}.
 \label{eq:alpha-cota-cola-uniforme}
\end{equation}
Moreover,
\begin{equation}
 \left|\frac{\mathrm d}{\mathrm dx}
  \left(\lambda\frac{x^{10}}{1-x^3/729}\right)\right|
 \leq 8.711\mathbin{\times}10^{-6}
 \left(\frac{10\,10^{-18}}{1-r}
 +\frac{3\,10^{-24}/729}{(1-r)^2}\right)
 <8.712\mathbin{\times}10^{-23}.
 \label{eq:alpha-cota-derivada-cola}
\end{equation}
The tail of the derivatives beyond block \(M\) also admits the
depth-dependent estimate
\begin{equation}
 \sup_{0\le x\le10^{-2}}
 \left|E'_{\mathsf s,\infty}(x)-E'_{\mathsf s,M}(x)\right|
 \le 8.711\mathbin{\times}10^{-24}\,r^{M+1}
 \left(\frac{13+3M}{1-r}+\frac{3r}{(1-r)^2}\right)
 \xrightarrow[M\to\infty]{}0.
 \label{eq:alpha-cota-resto-derivadas-rev08}
\end{equation}

\begin{theorem}[Limit function and unique simple root of the correlated chart]
\label{thm:alpha-raiz-completa-explicita}
The sequence \(E_{\mathsf s,M}\) converges uniformly, together with its
derivatives, to \(E_{\mathsf s,\infty}\) on
\(\overline I_\alpha=[0,10^{-2}]\). We have
\begin{equation}
 E_{\mathsf s,\infty}(a(\mathsf s))=0,
 \qquad
 \inf_{x\in\overline I_\alpha}
 E'_{\mathsf s,\infty}(x)>6.239.
 \label{eq:alpha-raiz-y-derivada-completas}
\end{equation}
Consequently, \(a(\mathsf s)\) is the unique simple root of the complete
function in \(I_\alpha\).
\end{theorem}

\begin{proof}
The ratio of two consecutive terms of the tail is bounded by \(r<1\),
which gives \eqref{eq:alpha-cota-cola-uniforme}; term-by-term
differentiation and summation of the two geometric series give
\eqref{eq:alpha-cota-derivada-cola} and
\eqref{eq:alpha-cota-resto-derivadas-rev08}. The proof for the jet establishes
\(E'_{9,\mathsf s}>6.24\) on \(\overline I_\alpha\); hence, the
complete perturbation leaves the derivative above \(6.239\).
Finally, substituting \eqref{eq:alpha-lambda-estado} into
\eqref{eq:alpha-funcion-completa-explicita}, the two summands evaluated at
\(a(\mathsf s)\) cancel exactly. The function is strictly
increasing; its unique root is therefore simple and coincides with that pre-coordinate.
\end{proof}

To make the identification effective at every depth, directed
endpoint evaluations first certify
\[
 E_{\mathsf s,\infty}(0)<0<
 E_{\mathsf s,\infty}(10^{-2}),
 \qquad
 E'_{\mathsf s,\infty}\geq\mu:=6239/1000>c:=6
 \quad\hbox{on }[0,10^{-2}].
\]
By continuity there is a root in the interval and, by the second inequality,
it is unique. Only after this existence, monotonicity and the invariant
that the bracket contains the root have been established is the mean-value-theorem
fallback applied. We start from \(D_0=[0,10^{-2}]\). If
\(D_j=[r_j,s_j]\), let \(m_j=(r_j+s_j)/2\), and let
\([F_j^-,F_j^+]\) be a directed rational enclosure of
\(E_{\mathsf s,\infty}(m_j)\), obtained by propagating the cylinders of
\(\pi_{\mathrm{HMT}},e_{\mathrm{HMT}},\varphi_{\mathrm{HMT}}\),
\(\delta\) and \(\lambda\), with
\begin{equation}
 0\leq F_j^+-F_j^-\leq \frac{c}{4}\,(s_j-r_j).
 \label{eq:alpha-anchura-evaluacion-rescate-rev08}
\end{equation}
When \(F_j^+<0\), the right half is retained;
when \(F_j^->0\), the left half is retained. If
\(F_j^-\leq0\leq F_j^+\), no attempt is made to decide whether the midpoint is
exactly the root. Let \(w_j=s_j-r_j\). The uniform bound
\(E'_{\mathsf s,\infty}\geq c=6\) and the mean value theorem imply that
any root contained in \(D_j\) belongs to
\begin{equation}
 D_{j+1}=D_j\cap
 \left[m_j-\frac{w_j}{4},
             m_j+\frac{w_j}{4}\right].
 \label{eq:alpha-rescate-derivada-rev08}
\end{equation}
Indeed, let \(\xi_j\) be the root and let
\(y_j=E_{\mathsf s,\infty}(m_j)\in[F_j^-,F_j^+]\). Since the enclosure contains
zero and satisfies \eqref{eq:alpha-anchura-evaluacion-rescate-rev08}, we have
\(|y_j|\leq cw_j/4\). The mean value theorem gives
\(|m_j-\xi_j|\leq |y_j|/c\leq w_j/4\), proving
\eqref{eq:alpha-rescate-derivada-rev08}. The argument includes the case
\(y_j=0\): the root is then at the midpoint, but the algorithm does not
need to recognize this equality in order to retain it.

If the available enclosure does not satisfy
\eqref{eq:alpha-anchura-evaluacion-rescate-rev08}, the source cylinders
are first refined and the evaluation is repeated; no sign is declared and no
insufficient contraction is accepted. In each of the three branches the new
interval retains the root and has diameter at most \(w_j/2\). For each
finite target depth, directed refinement allows the width
bound to be met without an uncertified comparison of a real number with zero.
Monotonicity proves inductively that
\begin{equation}
 a(\mathsf s)\in D_{j+1}\subseteq D_j,\qquad
 \operatorname{diam}D_{j+1}\leq
 \tfrac12\operatorname{diam}D_j,
 \qquad \operatorname{diam}D_j\longrightarrow0,
 \label{eq:alpha-biseccion-racional}
\end{equation}
where the inequality is a contraction bound, not an imposed equality.
The executable certificate applies
this rule to twenty-four base-one-thousand levels and checks that the complete
outer closure \(A_{1000}\), not only its lower endpoint, lies inside
each expressed cylinder. Its negative tests reject an excessively wide
enclosure, an incorrectly ordered enclosure, a point other than the midpoint, a domain without
a sign change, and a degenerate bracket. Coinductive extension allows
one thousand to be replaced by any required depth.

Specifically, for \(S_N=1000^N\) the verifier computes
\begin{equation}
 j_N^-:=\lfloor S_Na_-\rfloor,
 \qquad j_N^+:=\lfloor S_Na_+\rfloor
 \label{eq:alpha-indices-cilindro-dirigidos-rev08}
\end{equation}
and expresses the level only if \(j_N^-=j_N^+\). This equality, which includes the
upper endpoint of the outer closure, proves
\(A_{1000}\subset[j_N^-/S_N,(j_N^-+1)/S_N)\) and avoids assigning to the left
cell an endpoint lying exactly on its right boundary.

Let \(C_N^A(\mathsf s)\) be the integer cylinder in
\eqref{eq:alpha-cilindros-via-a-exactos}. We recursively choose
\(m(N)>m(N-1)\) so that
\(\operatorname{diam}D_{m(N)}<1000^{-N}/4\), and define
\begin{equation}
I_{B,N}:=D_{m(N)}\cap C_N^A(\mathsf s).
 \label{eq:alpha-intervalos-b-completos}
\end{equation}
Define
\(\ell_N:=\inf I_{B,N}\) and \(u_N:=\sup I_{B,N}\). Both are rational;
the interval is nonempty because both factors contain
\(a(\mathsf s)\). Since the cylinder is half-open, \(u_N\) may not
belong to \(I_{B,N}\). Evaluation at that endpoint is understood through
the continuous extension of \(E_{\mathsf s,\infty}\) to the outer closure. Thus,
\begin{equation}
 I_{B,N+1}\subseteq I_{B,N}\subseteq C_N^A(\mathsf s),
 \qquad \operatorname{diam}I_{B,N}\longrightarrow0,
 \qquad E_{\mathsf s,\infty}(\ell_N)\leq0\leq
 E_{\mathsf s,\infty}(u_N).
 \label{eq:alpha-inclusion-signos-completa}
\end{equation}
The bound~\eqref{eq:alpha-cota-cola-uniforme} allows a
truncation to be chosen that preserves the sign at each endpoint where the complete function
does not vanish. If an endpoint coincides with the root, the identity of
the complete function and the membership already proved are used; that same
sign is not attributed to its truncations. Indeed, the exact remainder is
\[
 E_{\mathsf s,\infty}(x)-E_{\mathsf s,M}(x)
 =\lambda(\mathsf s)x^{10}
   \frac{(qx^3)^{M+1}}{1-qx^3}.
\]
In the positive domain under consideration, where \(\lambda(\mathsf s)<0\), we
obtain at the root
\[
 E_{\mathsf s,M}(a(\mathsf s))
 =-\lambda(\mathsf s)a(\mathsf s)^{10}
   \frac{(qa(\mathsf s)^3)^{M+1}}{1-qa(\mathsf s)^3}>0
 \qquad(M<\infty).
\]
Therefore, the certification of the weak signs in
\eqref{eq:alpha-inclusion-signos-completa} refers to
\(E_{\mathsf s,\infty}\). Its evaluation uses the complete function or the
truncation accompanied by the signed remainder. This covers both the included
infimum and the supremum that may be excluded from the cylinder.

\begin{proposition}[Compatibility squares and common cylinders]
\label{prop:alpha-cuadrados-compatibilidad}
The restrictions of words and of analytic representations commute:
\[
\begin{array}{ccc}
 \mathcal X_\alpha^{\mathrm{enr}}&
 \xrightarrow{\ \mathsf T_{\alpha,N'}\ }&\mathcal W_{N'}\\[1mm]
 \Big\Vert&&\Big\downarrow\rho_{N',N}\\[1mm]
 \mathcal X_\alpha^{\mathrm{enr}}&
 \xrightarrow{\ \mathsf T_{\alpha,N}\ }&\mathcal W_N
\end{array}
\qquad
\begin{array}{ccc}
 \mathcal X_\alpha^{\mathrm{enr}}&
 \xrightarrow{\ \Gamma_{M'}\ }&\mathfrak J_{M'}\\[1mm]
 \Big\Vert&&\Big\downarrow\tau_{M',M}\\[1mm]
 \mathcal X_\alpha^{\mathrm{enr}}&
 \xrightarrow{\ \Gamma_M\ }&\mathfrak J_M .
\end{array}
\]
The collection \(I_{B,N}\) provides the effective connection between the two
squares and satisfies concretely
\(I_{B,N}\subseteq C_N^A\) for every \(N\).
\end{proposition}

\begin{proof}
The first square is the compatibility of integer normalization; the second
is Proposition~\ref{prop:alpha-naturalidad-recurrencia-explicita}. The
inclusion follows from the definition~\eqref{eq:alpha-intervalos-b-completos},
and nonemptiness and the signs follow from
Theorem~\ref{thm:alpha-raiz-completa-explicita}. No step identifies the
root of a truncation with that of an arbitrary tail.
\end{proof}

We now define, without leaving the prefix reader implicit,
\begin{equation}
 \alpha_B(\mathsf s):=
 \operatorname*{arg\,zero}_{x\in I_\alpha}E_{\mathsf s,\infty}(x),
 \qquad
 \operatorname{pref}_N(x):=
 \frac{\lfloor1000^Nx\rfloor}{1000^N},
 \qquad
 \alpha_{B,N}:=\operatorname{pref}_N(\alpha_B).
 \label{eq:alpha-definicion-b-y-prefijos-rev08}
\end{equation}

\begin{theorem}[Projective equality of the two correlated representations]
\label{thm:alpha-dos-publicaciones-completas}
For every \(N\), the integer and analytic routes express the same prefix:
\begin{equation}
 \boxed{\alpha_{A,N}=\operatorname{pref}_N(\alpha_A)
 =\operatorname{pref}_N(\alpha_B)=\alpha_{B,N}.}
 \label{eq:alpha-dos-vias-cada-nivel}
\end{equation}
In the limit,
\begin{equation}
 \boxed{\alpha_A
 =\operatorname*{arg\,zero}_{x\in I_\alpha}
       E_{\mathsf s,\infty}(x)
 =\alpha_B=\alpha_{\mathrm{HMT}}.}
 \label{eq:alpha-dos-vias-limite-coinductivo}
\end{equation}
\end{theorem}

\begin{proof}
Integer normalization expresses exactly \(a(\mathsf s)\).
Theorem~\ref{thm:alpha-raiz-completa-explicita} proves that the analytic
chart represents that same point as its unique simple root. The
half-open cylinders \(C_N^A\) fix a canonical prefix, and
\(I_{B,N}\subseteq C_N^A\) contains \(\alpha_B=\alpha_A\). Therefore, their
prefixes coincide for each \(N\); nesting and diameters tending to
zero give equality of the limits. Metrological
recognition occurs afterwards.
\end{proof}

The analytic function constitutes a correlated representation of the same
state, not a second derivation causally independent of the integer
closure. This precision does not weaken the equality proved: it identifies its
exact mechanism, prevents attributing to the jet a tail that it does not determine, and avoids
promoting a conventional value to a generator.


